\chapter*{致读者}
\phantomsection
\addcontentsline{toc}{chapter}{致读者}

\begin{enumerate}
\def\labelenumi{\arabic{enumi}.}
\item
  数学原理丛书从最开端讲起数学，并给出完整的证明。原则上，它不要求读者具备任何特别的数学知识，只要求对数学推理有一定的熟悉程度，以及对抽象思维有一定的能力。尽管如此，它特别适合于至少已较好地掌握大学数学课程头一两年内容的读者。
\item
  我们所选的阐述方法是公理化的，并且通常从一般过渡到特殊。证明的要求使题材具有严格固定的顺序。因此，除非读者已经具有相当广博的数学知识，否则某些考虑的用处可能要到较后的章节才会显现出来。
\item
  丛书分为若干卷（卷），每卷又分为若干章。法文版中已全部或部分出版的各卷列在下面。当有英译本可用时，相应的英文标题在括号中注明。在本卷中，引用一律指英译本（如果有的话），否则指法文版。
\end{enumerate}

Théorie des Ensembles（集合论）

记为

E

(Set Theory)

Algèbre（代数） \(^{1}\)

---

A

(Alg)

Topologie Générale（一般拓扑）

---

TG

(Gen.~Top.)

Fonctions d\textquotesingle une Variable Réelle（实变函数） \(^{2}\)

---

FVR

(FRV)

Espaces Vectoriels Topologiques（拓扑向量空间）

---

EVT

(Top. Vect. Sp.)

Intégration（积分）

---

INT

Algèbre Commutative（交换代数） \(^{3}\)

---

AC

(Comm. Alg.)

Variétés Différentielles et Analytiques（微分簇与解析簇）

---

VAR

Groupes et Algèbres de Lie（李群与李代数） \(^{4}\)

---

LIE

(LIE)

Théories Spectrales（谱理论）

TS

在前六卷（按上列次序）中，正文中的每个论断都只把那些已在同一章中、或在按如下次序排列的前面诸章中讨论过的结果作为已知：E；A，第 I 至 III 章；TG，第 I 至 III 章；A，从第 IV 章起；TG，从第 IV 章起；FVR；EVT；INT。

从第七卷起，读者通常可以找到关于该卷与其他各卷之间逻辑关系的明确说明（前六卷总假定已知）。

\begin{enumerate}
\def\labelenumi{\arabic{enumi}.}
\setcounter{enumi}{3}
\item
  不过，我们有时在正文中插入了这样的例子，它们涉及读者可能已经知道、但本丛书尚未讨论的事实。这类例子置于两个星号之间：*\ldots*。大多数读者无疑会发现这些例子有助于理解正文。在另一些情形，*\ldots* 之间的段落涉及本丛书其他地方所讨论的结果。我们希望读者能够验证其中不存在任何恶性循环。
\item
  每一章的逻辑框架由该章的定义、公理和定理组成。这些部分是今后使用时主要应当记住的。较次要的结果以及容易由定理推出的结果，标为“命题”、“引理”、“推论”、“注”等。初读时可以略去的部分以小字排印。对特别重要的定理的评注偶尔以“按语（scholium）”为名出现。
\end{enumerate}

为避免令人生厌的重复，有时方便的做法是引进只在某一章或某一章的某一节内有效的记号或缩写（例如，在只讨论交换环的一章中，“环”一词总是指“交换环”）。这类约定总是明确地予以说明，一般出现在采用它的那一章的开头。

\begin{enumerate}
\def\labelenumi{\arabic{enumi}.}
\setcounter{enumi}{5}
\item
  有些段落旨在预先提醒读者警惕严重的错误。这些段落在页边用记号（“危险弯道”）标出。
\item
  习题的目的，既在于使读者能确认自己已消化正文，也在于使读者注意到那些在正文中没有位置、但仍有意义的结果。最难的习题带有记号 ⌘。
\item
  一般说来，我们沿用通行的术语，除非看来有充分的理由偏离它。
\item
  我们特别力求始终使用严格正确的语言，而不牺牲简洁性。我们尽可能在正文中指出滥用语言之处；没有这类滥用，任何数学文本都有流于迂腐的危险，更谈不上可读性。
\item
  由于正文原则上是对理论作教条式的阐述，它一般不含对文献的引用。书目性资料汇集于历史注记之中。每个历史注记后面的参考书目一般只包含对所论理论的发展最为重要的那些书籍和原始论文。它绝不以完备自诩。
\end{enumerate}

至于习题，我们一般认为不值得注明其出处，因为它们取自许多不同的来源（原始论文、教科书、习题集）。

\begin{enumerate}
\def\labelenumi{\arabic{enumi}.}
\setcounter{enumi}{10}
\tightlist
\item
  在本卷中，对定理、公理、定义、\ldots{} 的引用依次按下列方式给出：
\end{enumerate}

-- 书（用第 3 节所列的缩写）、章和页码，即所说结果所在之处；

-- 仅当指本卷时，只给出章和页码。

结果综述用字母 R 来引用；因此 Set Theory, R 表示“集合论的结果综述”。
